\section*{Capítulo \ref{ch:g8:hormonal-communication} --- La comunicación hormonal}

\begin{solution}{exo:g8:hormonal-communication:1}
Que suelte su producto a la \omterm{prop:g4:journey-of-food:absorb}{sangre} y no por un conducto. La \omterm{def:g8:hormonal-communication:glands}{glándula} de
mando de la base del \omterm{def:g5:brain-in-command:system}{cerebro}: directora de las demás y la que arranca
la \omterm{def:g8:puberty:puberty}{pubertad}. El tiroides: el marcapasos del cuerpo. Las suprarrenales:
la \omterm{def:g8:hormonal-communication:glands}{hormona} de alarma. El páncreas: la gestión de la \omterm{prop:g7:organs-at-work:bill}{glucosa} de la
\omterm{prop:g4:journey-of-food:absorb}{sangre}. (\omterm{def:g8:reproductive-systems:female}{Ovarios} y \omterm{def:g8:reproductive-systems:male}{testículos}: las \omterm{def:g8:hormonal-communication:glands}{hormonas} del ciclo y de la
reconstrucción.)
\end{solution}

\begin{solution}{exo:g8:hormonal-communication:2}
El transporte es una emisión y la dirección se pone del lado del
receptor: solo responden las \omterm{def:g5:first-look-at-cells:cell}{células} equipadas con los receptores de
esa \omterm{def:g8:hormonal-communication:glands}{hormona}. La carta le llega a todas las casas; solo la abre la llave
que encaja.
\end{solution}

\begin{solution}{exo:g8:hormonal-communication:3}
El ritmo del \omterm{def:g4:heart-and-blood:heart}{corazón} sube (\cref{ch:g7:heart-and-circulation}); las
vías respiratorias se ensanchan
(\cref{ch:g7:lungs-and-gas-exchange}); se suelta el banco de \omterm{prop:g7:organs-at-work:bill}{glucosa}
del \omterm{prop:g7:digestion-and-nutrients:absorption}{hígado} (\cref{ch:g7:digestion-and-nutrients}); la \omterm{prop:g4:journey-of-food:absorb}{sangre} se desvía
a los \omterm{def:g3:bones-and-muscles:muscle}{músculos} (\cref{ch:g7:organs-at-work}) --- y la \omterm{def:g1:the-five-senses:sense}{piel} pálida y con
hormigueo completa el juego.
\end{solution}

\begin{solution}{exo:g8:hormonal-communication:4}
La \omterm{def:g8:hormonal-communication:glands}{glándula} que escribe vigila el resultado de sus cartas y frena a
medida que se alcanza el objetivo --- termostato, no grifo. Ejemplo
estrella: la \omterm{prop:g7:organs-at-work:bill}{glucosa} de la \omterm{prop:g4:journey-of-food:absorb}{sangre}, ingresada al subir y soltada al
bajar por las \omterm{def:g8:hormonal-communication:glands}{hormonas} emparejadas del páncreas.
\end{solution}

\begin{solution}{exo:g8:hormonal-communication:5}
La señal de la \omterm{def:g8:hormonal-communication:glands}{hormona} de almacenamiento --- sin escribir, o sin que la
respondan sus receptores. La \omterm{prop:g7:organs-at-work:bill}{glucosa} se amontona entonces después de
las comidas, desborda la \omterm{def:g7:organs-at-work:recovery}{recuperación} del \omterm{def:g7:removing-wastes:kidneys}{riñón} y aparece en la \omterm{def:g7:removing-wastes:kidneys}{orina}.
\end{solution}

\begin{solution}{exo:g8:hormonal-communication:6}
En la base del \omterm{def:g5:brain-in-command:system}{cerebro} --- la \omterm{def:g8:hormonal-communication:glands}{glándula} de mando, con el \omterm{def:g5:brain-in-command:system}{sistema
nervioso} por encima y las \omterm{def:g8:hormonal-communication:glands}{hormonas} por debajo. Ejemplo: el susto ---
primero las señales nerviosas y después la carta de las suprarrenales
sosteniendo la alerta; o la duración del día convirtiéndose en las
\omterm{def:g8:hormonal-communication:glands}{hormonas} de cría de las ovejas.
\end{solution}

\begin{solution}{exo:g8:hormonal-communication:7}
Un grifo echa hasta que alguien lo cierra; un termostato mide lo que
está cambiando y se ajusta solo. La \omterm{def:g8:hormonal-communication:glands}{glándula} lee el nivel que regula
--- la secreción sube y se apaga con la necesidad, sin que haga falta
ninguna mano de fuera.
\end{solution}

\begin{solution}{exo:g8:hormonal-communication:8}
Escritora: el páncreas, disparado por la \omterm{prop:g7:organs-at-work:bill}{glucosa} de la \omterm{prop:g4:journey-of-food:absorb}{sangre} que sube.
Receptores: el \omterm{prop:g7:digestion-and-nutrients:absorption}{hígado} y los \omterm{def:g3:bones-and-muscles:muscle}{músculos} sobre todo --- ingresan el
excedente cuando les llega. \omterm{prop:g8:hormonal-communication:feedback}{Retroalimentación}: el propio nivel que baja
apaga la secreción. Fallo y uso: la \omterm{ex:g8:hormonal-communication:diabetes}{diabetes} --- y las dosis medidas de
la medicina, la carta por jeringuilla.
\end{solution}

\begin{solution}{exo:g8:hormonal-communication:9}
Las dos escriben a mano una carta del cuerpo: la jeringuilla aporta el
mensaje de almacenamiento que le falta al páncreas en las horas de
comer; la píldora aporta a diario el mensaje de ``suspendido'' del
ciclo. El mismo acto --- imitar una \omterm{def:g8:hormonal-communication:glands}{hormona} para llevar su sistema a
propósito --- una restaurando una correspondencia que ha fallado y la
otra imponiéndose por elección a una que está sana.
\end{solution}

\begin{solution}{exo:g8:hormonal-communication:10}
El arranque de los tiritones: nervioso --- órdenes musculares en una
fracción de segundo. La barba: hormonal --- meses, todo el cuerpo,
torrente sanguíneo. El carril de la \omterm{prop:g7:organs-at-work:bill}{glucosa}: hormonal ---
correspondencia con \omterm{prop:g8:hormonal-communication:feedback}{retroalimentación}. El reflejo rotuliano: nervioso
--- un empalme, una cincuentava parte de segundo.
\end{solution}

\begin{solution}{exo:g8:hormonal-communication:11}
La percepción por los \omterm{def:g1:the-five-senses:sense}{ojos} de los días que se acortan es nerviosa; la
oficina de la base del \omterm{def:g5:brain-in-command:system}{cerebro} la traduce y dirige a la \omterm{def:g8:hormonal-communication:glands}{glándula} de
mando; las cartas de esta despiertan el programa hormonal propio de los
\omterm{def:g8:reproductive-systems:female}{ovarios}; y arrancan los ciclos de las ovejas --- entra \omterm{def:g6:exploring-our-environment:conditions}{luz} por el
telégrafo y salen corderos por correo.
\end{solution}

\begin{solution}{exo:g8:hormonal-communication:12}
Defensa: una sola \omterm{prop:g4:journey-of-food:absorb}{sangre} lleva todas las cartas y, aun así, cada una
actúa solo donde responden los receptores --- el público define el
canal, así que la emisión se convierte en líneas privadas. La
prolongación: una \omterm{def:g5:first-look-at-cells:cell}{célula} se une a un público construyendo el receptor
--- así que qué tejidos escuchan se construye también, órgano a órgano,
durante el desarrollo; las dianas de la \omterm{def:g8:puberty:puberty}{pubertad} --- las cuerdas
vocales, la \omterm{def:g1:the-five-senses:sense}{piel}, los cartílagos de crecimiento --- estaban equipadas
con sus receptores por adelantado, esperando la primera carta durante
años.
\end{solution}

\begin{solution}{pb:g8:hormonal-communication:1}
\textbf{1.} Las cartas emparejadas del páncreas --- la \omterm{def:g8:hormonal-communication:glands}{hormona} de
almacenamiento al subir y su compañera que suelta al bajar ---
manteniendo el carril por \omterm{prop:g8:hormonal-communication:feedback}{retroalimentación} mientras su dueña dormía.

\textbf{2.} Disparador: la \omterm{prop:g7:organs-at-work:bill}{glucosa} de la comida atravesando las
\omterm{prop:g7:digestion-and-nutrients:absorption}{vellosidades} y subiendo el nivel de la \omterm{prop:g4:journey-of-food:absorb}{sangre}. Escritora: el páncreas.
Carta: la \omterm{def:g8:hormonal-communication:glands}{hormona} de almacenamiento. Receptores: el \omterm{prop:g7:digestion-and-nutrients:absorption}{hígado} y los
\omterm{def:g3:bones-and-muscles:muscle}{músculos}, que ingresan el excedente.

\textbf{3.} Un grifo se pasaría del objetivo hasta un hundimiento; que
el descenso se nivele dentro del carril enseña que la secreción se
apaga a medida que llega su propio resultado --- se mide sola y se para
sola.

\textbf{4.} Que subió con el nivel y se murió con él: una curva de
secreción a la sombra de la curva de la \omterm{prop:g7:organs-at-work:bill}{glucosa}, con el máximo
siguiendo al máximo --- el recorrido del termostato.

\textbf{5.} La compañera que suelta (con la ayuda de la \omterm{def:g8:hormonal-communication:glands}{hormona} de
alarma en el esfuerzo duro): la \omterm{prop:g7:organs-at-work:bill}{glucosa} ingresada del \omterm{prop:g7:digestion-and-nutrients:absorption}{hígado} --- el
banco llenado en el desayuno a la manera del
\cref{ex:g7:digestion-and-nutrients:breadtrace} --- gastada en la
factura de los \omterm{def:g3:bones-and-muscles:muscle}{músculos} que trabajan.

\textbf{6.} Esa misma compañera que suelta: abriendo el banco de manera
constante durante las horas sin comida y ajustándose al gasto del
cuerpo --- el carril mantenido desde el otro lado.

\textbf{7.} Las \omterm{def:g8:hormonal-communication:glands}{glándulas} suprarrenales. El azúcar prepara combustible
para \emph{huir o pelear}: el sentido entero de la carta de alarma es
un cuerpo abastecido y desviado para un esfuerzo inmediato --- llegue
ese esfuerzo o no.

\textbf{8.} Porque la correspondencia sana le respondió: el excedente
soltado y no gastado (el susto pasó) disparó a su vez la carta de
almacenamiento, y el carril se restauró --- alarma y termostato uno
detrás del otro.

\textbf{9.} Sin tratamiento: la subida del desayuno siguiendo hacia
arriba y sin descender --- ninguna carta de almacenamiento escrita ni
oída --- una meseta alta durante toda la mañana; y pasado el umbral de
la \omterm{def:g7:organs-at-work:recovery}{recuperación}, aparece azúcar en la \omterm{def:g7:removing-wastes:kidneys}{orina} (la señal del
\cref{ex:g7:removing-wastes:thrift}).

\textbf{10.} La propia carta de las horas de comer del páncreas. La
dosis tiene que ajustarse al plato porque el tamaño de la carta tiene
que ajustarse a la \omterm{prop:g7:organs-at-work:bill}{glucosa} que va a ingresar --- si es pequeña deja un
amontonamiento y si es grande se pasa de carril hacia abajo, que es una
urgencia por sí misma.

\textbf{11.} Noche y tarde: la pareja del páncreas, con el \omterm{prop:g7:digestion-and-nutrients:absorption}{hígado} y los
\omterm{def:g3:bones-and-muscles:muscle}{músculos} respondiendo --- carril mantenido. Comidas de la mañana y de
la noche: habla la \omterm{prop:g7:organs-at-work:bill}{glucosa}, escribe el páncreas, ingresa el \omterm{prop:g7:digestion-and-nutrients:absorption}{hígado}.
Mediodía: gasta el esfuerzo y la compañera que suelta abre el banco.
Las cinco: interrumpen las suprarrenales y el páncreas recoge después.
Un solo valor, cinco hablantes y ni un silencio en todo el día.

\textbf{12.} ``Un sistema de \omterm{prop:g8:hormonal-communication:feedback}{retroalimentación}: escritoras que leen los
resultados de sus propias cartas --- que es como un valor circula por
un carril estrecho toda una vida sin nadie al volante''.
\end{solution}
